\documentclass[11pt]{article}

\usepackage[margin=1in]{geometry}
\usepackage{amsmath,amssymb,amsthm,mathtools}
\usepackage{enumitem}
\usepackage{hyperref}
\hypersetup{colorlinks=true,linkcolor=blue,citecolor=blue,urlcolor=blue}
\usepackage{orcidlink}
\usepackage{fontawesome5}

\newtheorem{theorem}{Theorem}[section]
\newtheorem{proposition}[theorem]{Proposition}
\newtheorem{lemma}[theorem]{Lemma}
\newtheorem{corollary}[theorem]{Corollary}
\newtheorem{empirical}[theorem]{Empirical Result}
\theoremstyle{remark}
\newtheorem{remark}[theorem]{Remark}

\title{Wirsching's 2003 Conjectures on the Base-$3$ Fabius Function:\\
A Proof, a Microcanonical Decomposition, and a Certified Numerical Test}

\author{Renato Augusto Tavares~{\small\href{mailto:dr.renatotavares@gmail.com}{\textcolor{black}{\faEnvelope}}}~\orcidlink{0009-0002-0196-3311}\\
\normalsize Universidade Federal de Goiás\\
\small \href{https://orcid.org/0009-0002-0196-3311}{https://orcid.org/0009-0002-0196-3311}}
\date{\today}

\begin{document}

\maketitle

\begin{abstract}
% DRAFT -- rewrite by hand before submission (Rule 4b: abstracts carry
% the highest marker density of any section in the paper).
Wirsching (2003) reduces uniform positive predecessor density for the
$3n+1$ map to a chain of five conditions, organized into three
conjectures. Conjecture 1 concerns Wirsching's path-counting
generators. Conjecture 2 concerns a one-dimensional averaging operator
$W_3$, whose unique $L^1$ fixed point $\varphi$ is a base-$3$ analogue
of the Fabius function. Conjecture 3 concerns the asymptotics of
$\varphi$ itself. We prove Conjecture 1 by a generating-function
cancellation. Reading the primary source directly, we find that
Conjecture 2's hypothesis carries no information about the residue
coordinate the rest of the chain needs: the only bridge the source
supplies is a strong-convergence theorem, uniform on bounded families
but not in growing resolution. We give a microcanonical decomposition
of the Syracuse law in terms of Wirsching's generators, prove an exact
equivalence-of-ensembles theorem at fixed and sublinear precision, and
prove that the same decomposition's latent cost vectors fail to
converge at linear precision, a narrower fact than nonequivalence
after projection, which remains open. A certified evaluation of
$\varphi$ to error $10^{-8}$ through depth $\ell=500$ supports
Conjecture 3, subject to a wider $\pm0.015$ extrapolation uncertainty
outside that certificate.
\end{abstract}

\section{Introduction}
\label{sec:intro}

Wirsching~\cite{Wirsching2003} studies uniform positive predecessor
density for the Collatz map through a functional equation for a
specific probability density $\varphi$, a base-$3$ analogue of the
Fabius function, arising as the limiting density of a normalized path
count in the Collatz reverse tree. The paper reduces the target
density statement to a chain of implications through three
increasingly concrete conjectures, of which only the first has since
been settled; the third has numerical support in this line of work,
and the second, on which the whole chain pivots, has remained open.

\S\ref{sec:chain} fixes Wirsching's objects and states the chain.
\S\ref{sec:conj1} proves the first conjecture.
\S\ref{sec:microcanonical} develops the microcanonical decomposition of
the Syracuse law that connects Wirsching's path counts to the weighted
covering formulation of the endogeny-barrier literature this line of
work also produces \cite{BarrierCompanion}, together with the
equivalence- and nonequivalence-of-ensembles theorems that bound how
much of the second conjecture's target local-limit methods reach.
\S\ref{sec:conj3} reports the certified numerical test of the third
conjecture.

\section{Wirsching's chain, from the primary source}
\label{sec:chain}

Wirsching's central combinatorial objects are the ``Elka'' functions
$e_\ell(k,a):=|E_{\ell,k}(a)|$, where $E_{\ell,k}(a)$ is the set of
paths $b\xrightarrow{T}\cdots\xrightarrow{T}a$ in the Collatz graph
with $T=T_0$ exactly $k$ times and $T=T_1$ exactly $\ell$ times
\cite[\S1]{Wirsching2003} ($T_0,T_1$ the two branches of the
accelerated $3n+1$ map), viewed on the state space
$\mathbb X=\mathbb I_3\times\mathbb Z_3^\times$: the
$\mathbb Z_3^\times$ factor carries the $3$-adic variable habitat of
the Syracuse measure, while the $\mathbb I_3\cong\{0,1,2\}^{\mathbb
N}$ factor carries the normalized step count $k/3^\ell$, converging to
a density $\varphi$.

The compactly supported \emph{generators} $g_\ell(k,a)$ are defined by
the recursion \cite[\S2, eq.~(2.1)]{Wirsching2003}
\[
 g_{\ell+1}(k,a)=\sum_{0\le j<2\cdot3^\ell}
 g_\ell\bigl(k-j,\tfrac{2^{j+1}a-1}{3}\bigr),
 \qquad g_0(k,a)=\mathbf 1[k=0]\,\mathbf1[a\in\mathbb Z_3],
\]
linked to the Elka functions by $e_\ell=p_\ell*g_\ell$, where
$p_\ell(m)$ counts the ways to pay $m$ with coins $c_0=1$,
$c_j=2\cdot3^{j-1}$ ($j\ge1$) \cite[\S2]{Wirsching2003}. The class
$A_\delta$ consists of the integer sequences $(k_\ell)$ with
$|\ell-k_\ell|\le\delta\sqrt\ell$ for every $\ell$
\cite[eq.~(1.5)]{Wirsching2003}; $\widetilde A_\delta$ is the
analogous class of sequences $(x_\ell)\in\mathbb I_3^{\mathbb N}$ used
once the objects are transported to the continuous state space of
\S\ref{sec:microcanonical} below. $\varphi_0$ is the explicit
asymptotic of Berg and Krüppel for the solution of a truncated
functional equation satisfied by $\varphi$
\cite[eq.~(7.11)]{Wirsching2003}, used in Conjecture~3 below in place
of $\varphi$ itself, which has no known closed form.

\begin{proposition}[Base-$3$ Fabius function; Wirsching 2003, Corollaries 7 and 8]\label{prop:fabius}
$\varphi$ is the density of $X=\sum_{j\ge1}2U_j\cdot3^{-j}$, $U_j$
i.i.d.\ uniform on $[0,1]$, the base-$3$ analogue of the Fabius
function (which satisfies $f'(x)=2f(2x)$; here $\varphi'(x)=
\tfrac92\varphi(3x)$ on $[0,2/3]$ by \cite[eq.~(7.7)]{Wirsching2003}).
It is the unique $L^1([0,1])$ fixed point of
the averaging operator $W_3f(x):=\tfrac32\int_{3x-2}^{3x}f(t)\,dt$,
$C^\infty$, piecewise polynomial away from the standard Cantor set
\cite[Corollary~7]{Wirsching2003}, with certified geometric convergence
$\|f_n-\varphi\|_1\le2^{-n+1}\|f_1-f_0\|_1$
\cite[Corollary~8]{Wirsching2003}.
\end{proposition}

Wirsching's paper reduces uniform positive predecessor density (that
the sets of predecessors of $T$ have density bounded below uniformly
over admissible targets) to a chain of five conditions
\cite[Definition~1, \S\S1--7]{Wirsching2003}:
\[
\begin{array}{ll}
\text{Theorem 1 (proved in \cite{Wirsching2003})} & (\star1)\Rightarrow
  \text{Positive Density}\\
\text{Conjecture 1} & (\star2)\Rightarrow(\star1)\\
\text{Theorem 2 (proved in \cite{Wirsching2003})} & (\star3)\Rightarrow(\star2)\\
\text{Conjecture 2} & (\star4)\Rightarrow(\star3)\\
\text{Conjecture 3} & (\star5);\ (\star5)\Rightarrow(\star4)\text{ proved in
  \cite[\S7]{Wirsching2003}}
\end{array}
\]
where $(\star1)$ is a $\liminf$ lower bound on Wirsching's Elka functions
relative to their Haar average, uniform over sequences
$(k_\ell)$ with $|\ell-k_\ell|\le\delta\sqrt\ell$; $(\star2)$ is the same
statement for his generators $g_\ell(k,a)$; $(\star3)$ upgrades the
$\liminf$ over $\ell$ to a bound holding for every $\ell\ge\ell_0$;
and $(\star4)$ is
\[
 \liminf_\ell\frac{(W_3^\ell\chi_1)(x_\ell)}{(W_3^\ell\chi_0)(x_\ell)}
 \ge\mu,\qquad\text{uniformly for }(x_\ell)\in\widetilde A_\delta,
\]
with $\chi_0=\mathbf 1_{[0,2/3]}$, $\chi_1=\mathbf1_{[1/3,1]}$, and
$W_3$ the averaging operator of Proposition~\ref{prop:fabius}. Three
facts about this chain matter for what follows.

First, $(\star4)$ is a statement purely about the one-dimensional operator
$W_3$: it never mentions Wirsching's generators, and carries no
information about the residue coordinate that $(\star3)$ needs. The
source offers two bridges between the two, not one: the averaged-generator
identity $\lim_\ell\int_{\mathbb Z_3^\times}\widetilde g_\ell(x,a)\,da=
\varphi(x)$ for fixed $x\in\mathbb I_3$ \cite[\S6]{Wirsching2003}, and
Wirsching's Theorem~3, strong convergence $S_\ell\to S_\infty$ on
equicontinuous bounded families, uniform in that family but not in
growing resolution. Neither is uniform in growing resolution, so
neither closes the gap between $(\star4)$ and $(\star3)$.

Second, $(\star3)$ and $(\star2)$ are the same inequality. Writing
$\widetilde g_\ell(x,a)=\gamma_\ell\,g_\ell(\lfloor3^\ell
x\rfloor,a)$ for a constant $\gamma_\ell$ not depending on $a$, the
constant cancels from both sides of the ratio each condition bounds;
what Wirsching's Theorem~2 adds is a change of quantifier, from a
$\liminf$ in $\ell$ at each $a$ to a bound holding at every
$\ell\ge\ell_0$ and every unit simultaneously.

Third, Theorem~1's own hypothesis and conclusion are about $a\in\mathbb
N$, non-cyclic, $a\not\equiv0\pmod3$: conditions $(\star2)$ and $(\star3)$,
stated over $\mathbb Z_3^\times$, are strictly stronger than what the
chain from Theorem~1 alone requires.

\section{Wirsching's Conjecture 1}
\label{sec:conj1}

\begin{theorem}[Wirsching's Conjecture 1]\label{thm:wirsching-conj1}
In the notation of \S\ref{sec:chain}, condition $(\star2)$ implies
condition $(\star1)$, with a possibly smaller central-limit window.
\end{theorem}

\begin{proof}
Put $c_0=1$ and $c_j=2\cdot3^{j-1}$ for $j\ge1$. If $q_\ell(k)$ counts
the bounded urn distributions defining $\bar g_\ell(k)$, then
\[
 \sum_{k\ge0}q_\ell(k)z^k
 =\prod_{j=0}^\ell\frac{1-z^{c_j}}{1-z},\qquad
 \sum_{m\ge0}p_\ell(m)z^m
 =\prod_{j=0}^\ell(1-z^{c_j})^{-1}.
\]
Since $q_\ell(k)=2\cdot3^{\ell-1}\bar g_\ell(k)$ and
$\bar e_\ell=p_\ell*\bar g_\ell$, cancellation gives
\[
 \bar e_\ell(k)=\frac{1}{2\cdot3^{\ell-1}}
                 \binom{k+\ell}{\ell}.
\]
The unrestricted partition count for the coins $1,2,6,18,\ldots$ is
at most $\exp(C\log^2(m+2))$. Also
\[
 \frac{\binom{k-m+\ell}{\ell}}{\binom{k+\ell}{\ell}}
 \le \left(\frac{k}{k+\ell}\right)^m.
\]
Uniformly for $|k-\ell|\le\delta\sqrt\ell$, the part of the convolution
with $m\ge\eta\sqrt\ell$ is
$\exp(-c\sqrt\ell+O(\log^2\ell))$. Choose
$\delta+\eta<\delta_1$. On the remaining terms $j=k-m$ lies in the
window of $(\star2)$, and positivity transfers its lower bound from
$g_\ell/\bar g_\ell$ to $e_\ell/\bar e_\ell$.
\end{proof}

\begin{remark}[Relation to the source]\label{rem:wirsching-conj1-source}
Wirsching records the two generating-function identities and the
closed form for $\bar e_\ell(k)$ used above \cite[\S\S1--2,
eq.~(1.4)]{Wirsching2003}, and the urn interpretation of
$q_\ell(k)=2\cdot3^{\ell-1}\bar g_\ell(k)$
\cite{Wirsching1998Urns}, and states the implication itself as
Conjecture~1 \cite[\S\S1--2]{Wirsching2003}. What the proof above adds
is the convolution-tail estimate and the window choice
$\delta+\eta<\delta_1$ that closes the implication, checked
coefficient by coefficient in the accompanying repository.
\end{remark}

The second conjecture remains open. The most concrete sufficient
condition above it is Conjecture~3, an asymptotic assertion on
$\varphi(z_\ell)/\varphi_0(z_\ell)$ along the central-limit window
$|\ell-k_\ell|\le\delta\sqrt\ell$, which we test numerically in
\S\ref{sec:conj3}.

\section{A microcanonical decomposition, and how far local-limit methods reach}
\label{sec:microcanonical}

\begin{theorem}[Microcanonical decomposition]\label{thm:microcanonical}
Let $g_\ell(k,a)$ be Wirsching's generator, let
$c_i=2\cdot3^i$ for $i\ge0$ (the coin $c_j$ of \S\ref{sec:conj1}
reindexed by $i=j-1$, dropping the extra coin $c_0=1$ that appears
there only), and let $\mu_\ell$ be the Syracuse law modulo
$3^\ell$. For every unit residue $a$,
\begin{equation}\label{eq:microcanonical}
 \mu_\ell(a)=
 \left(\prod_{i=0}^{\ell-1}\frac{1}{2(1-2^{-c_i})}\right)
 \sum_{k=0}^{3^\ell-\ell-1}2^{-k}g_\ell(k,a).
\end{equation}
Suppose there are $\delta,\eta>0$ such that, for all sufficiently large
$\ell$, every unit $a$, and every integer
$|k-\ell|\le\delta\sqrt\ell$,
\[
 g_\ell(k,a)\ge\eta\,\overline g_\ell(k),
 \qquad
 \overline g_\ell(k)=\frac{1}{2\cdot3^{\ell-1}}
 \sum_{b\in(\mathbb Z/3^\ell\mathbb Z)^\times}g_\ell(k,b).
\]
Then $\mu_\ell(a)\ge C_{\delta,\eta}3^{-\ell}$ uniformly in $a$ and
$\ell$. In particular, condition $(\star3)$ implies a uniform lower bound
on every cylinder of the Syracuse law, the estimate the weighted
covering formulation of the endogeny-barrier literature needs
\cite{BarrierCompanion}. The hypothesis $g_\ell(k,a)\ge\eta\bar
g_\ell(k)$ is open: Empirical Result~\ref{thm:microcanonical-finite}
below finds the minimum multiplicity is zero at $k=\ell$ at every
level computed, so no $\eta>0$ works there yet; the theorem is
conditional on closing that gap, not a report that it is closed.
\end{theorem}

\begin{proof}
At recursion level $i$, multiplication by $2$ has order
$c_i=2\cdot3^i$ modulo $3^{i+1}$. If $A_i$ has the geometric law
$\Pr(A_i=m)=2^{-m}$, then $J_i=A_i-1\bmod c_i$ satisfies
\[
 \Pr(J_i=j)=\frac{2^{-(j+1)}}{1-2^{-c_i}},
 \qquad 0\le j<c_i.
\]
Folding each exponent independently by the applicable order leaves the
Syracuse residue unchanged: from the recursion $a=(3a'+1)\cdot
2^{-(j+1)}$, the residue $a\bmod3^{i+1}$ is a function of $a'\bmod3^i$
and of $j\bmod c_i$ alone, since $c_i=2\cdot3^i$ is exactly the order
of $2$ modulo $3^{i+1}$. Wirsching's recursion for $g_\ell$ counts
the folded vectors with $\sum_iJ_i=k$. Multiplying their probabilities
gives \eqref{eq:microcanonical}.

Put $K_\ell=\sum_{i<\ell}J_i$. The summands are independent and are
uniformly dominated by an untruncated geometric variable. Direct
summation gives
$\mathbb E K_\ell=\ell+O(1)$ and
$\operatorname{Var}K_\ell=2\ell+O(1)$. The Lindeberg central limit
theorem therefore gives
\[
 \Pr(|K_\ell-\ell|\le\delta\sqrt\ell)
 \longrightarrow 2\Phi(\delta/\sqrt2)-1>0.
\]
Sum \eqref{eq:microcanonical} over the central values of $k$ and use
the assumed lower bound. Since there are $2\cdot3^{\ell-1}$ unit
residues, this yields
\[
 \mu_\ell(a)\ge
 \frac{\eta}{2\cdot3^{\ell-1}}
 \Pr(|K_\ell-\ell|\le\delta\sqrt\ell)
 \ge C_{\delta,\eta}3^{-\ell}.
\]
\end{proof}

\begin{proposition}[Conditional local-limit bridge]
\label{prop:complex-deconditioning}
Set
\[
 G_{\ell,a}(z)=\sum_k g_\ell(k,a)z^k,\qquad
 \overline G_\ell(z)=\frac{1}{2\cdot3^{\ell-1}}\sum_aG_{\ell,a}(z).
\]
Define the conditional cost laws
\[
 Q_{\ell,a}(k)=\frac{2^{-k}g_\ell(k,a)}{G_{\ell,a}(1/2)},
 \qquad
 \overline Q_\ell(k)=
 \frac{2^{-k}\overline g_\ell(k)}{\overline G_\ell(1/2)}.
\]
Then the exact identity
\begin{equation}\label{eq:deconditioning}
 \frac{g_\ell(k,a)}{\overline g_\ell(k)}
 =
 \frac{G_{\ell,a}(1/2)}{\overline G_\ell(1/2)}
 \frac{Q_{\ell,a}(k)}{\overline Q_\ell(k)}
\end{equation}
holds. Fix $\delta>0$. If, uniformly in unit $a$ and
$|k-\ell|\le\delta\sqrt\ell$,
\[
 \frac{G_{\ell,a}(1/2)}{\overline G_\ell(1/2)}\ge\eta>0,
 \qquad
 Q_{\ell,a}(k)\ge c_\delta\ell^{-1/2},
\]
then condition $(\star3)$ holds. Both hypotheses are open at $k=\ell$
for the same reason as in Theorem~\ref{thm:microcanonical}: the
minimum computed multiplicity is zero at every level tested.
\end{proposition}

\begin{proof}
Identity \eqref{eq:deconditioning} follows by cancellation. Under the
folded geometric law in the proof of Theorem~\ref{thm:microcanonical},
$\overline Q_\ell$ is the law of $K_\ell$.
The summands have uniformly bounded exponential moments,
mean sum $\ell+O(1)$, and variance $2\ell+O(1)$. Their lattice local
central limit theorem gives, uniformly in the stated window,
\[
 \overline Q_\ell(k)\asymp_\delta\ell^{-1/2}.
\]
The two assumed lower bounds and \eqref{eq:deconditioning} now give a
positive uniform lower bound for
$g_\ell(k,a)/\overline g_\ell(k)$.
\end{proof}

The first factor in \eqref{eq:deconditioning} is the canonical Syracuse
mass relative to its mean. The conjectural weighted-covering estimate
gives only a subexponential lower bound for this factor; the
proposition assumes a constant lower bound. The second factor is a
local limit theorem after conditioning on the residue. Its Fourier
representation is
\[
 Q_{\ell,a}(k)=\frac{1}{2\pi}\int_{-\pi}^{\pi}
 e^{-ikt}\frac{G_{\ell,a}(e^{it}/2)}{G_{\ell,a}(1/2)}\,dt.
\]

\begin{remark}[Local limit theorems for operator cocycles]
Spectral local limit theorems exist for random expanding dynamics
\cite{DragicevicEtAl2018} and for time-dependent expanding systems
\cite{Hafouta2020}. Their hypotheses include a twisted operator
cocycle acting on a fixed Banach space, a uniform spectral or
Ruelle--Perron--Frobenius gap near zero, and an aperiodicity estimate
away from zero. These hypotheses would give the phase control in
Proposition~\ref{prop:complex-deconditioning}. They have not been
verified for Wirsching's $S_\ell$. Theorem~3 of
\cite{Wirsching2003} gives strong convergence on fixed
equicontinuous families, while the densities used here acquire
resolution $3^{-\ell}$. The available local limit theorems do not
remove that growing-resolution obstruction.
\end{remark}

\begin{theorem}[Equivalence of ensembles at fixed precision]
\label{thm:fixed-precision-ensemble}
For a unit residue $a$ put
\[
 p_{\ell,k}(a)=\frac{g_\ell(k,a)}{\sum_b g_\ell(k,b)}.
\]
Fix $r\ge1$ and $\delta>0$. Uniformly over integers
$|k-\ell|\le\delta\sqrt\ell$, the projection of $p_{\ell,k}$ modulo
$3^r$ converges in total variation to the Syracuse law $\mu_r$ as
$\ell\to\infty$.
\end{theorem}

\begin{proof}
Under the canonical weights in Theorem~\ref{thm:microcanonical}, the
folded costs $J_0,\ldots,J_{\ell-1}$ are independent and
$p_{\ell,k}$ is the law of the resulting residue conditional on
$K_\ell=\sum_iJ_i=k$. Fix the terminal block of $r$ costs. For each
fixed value $\mathbf j$ of that block, with coordinate sum
$s(\mathbf j)$, independence gives
\[
 \Pr(\mathbf J=\mathbf j\mid K_\ell=k)
 =\Pr(\mathbf J=\mathbf j)
   \frac{\Pr(K_{\ell-r}=k-s(\mathbf j))}
        {\Pr(K_\ell=k)}.
\]
The lattice local central limit theorem used in
Proposition~\ref{prop:complex-deconditioning} makes the ratio tend to
one, uniformly in the stated window, for every fixed $\mathbf j$.
The geometric domination of the coordinates makes the block tails
uniformly summable. Splitting at block sum $o(\sqrt\ell)$ and applying
the same local estimate there extends the convergence to total
variation on the whole block.

The terminal $r$ folded costs determine the final residue modulo
$3^r$: pairing $J_{\ell-1-i}$, for $i=0,\ldots,r-1$, against the
reduced cap $2\cdot3^{r-1-i}$ is valid because the original cap
$c_{\ell-1-i}=2\cdot3^{\ell-1-i}$ is a multiple of $2\cdot3^{r-1-i}$
for every $r\le\ell$, which is exactly the divisibility $d\mid C$ used
below. At a terminal stage whose original cap is $C$, reduction to the
cap $d=2\cdot3^i$ required at precision $r$ is exact because $d\mid C$
and
\[
 \sum_{\substack{0\le j<C\\j\equiv s\ ({\rm mod}\ d)}}
 \frac{2^{-(j+1)}}{1-2^{-C}}
 =\frac{2^{-(s+1)}}{1-2^{-d}},\qquad 0\le s<d.
\]
The reduced block therefore has exactly the canonical folded-geometric
law defining $\mu_r$. Total variation cannot increase under the
deterministic map from the block to its residue, which proves the
claim.
\end{proof}

\begin{theorem}[Equivalence of ensembles at sublinear precision]
\label{thm:sublinear-precision-ensemble}
Let $r=r_\ell$ be an integer sequence with $1\le r\le\ell$ and
$r=o(\ell)$. Fix $\delta>0$. Uniformly over integers
$|k-\ell|\le\delta\sqrt\ell$, the projection of $p_{\ell,k}$ modulo
$3^r$ has total-variation distance $o(1)$ from $\mu_r$.
\end{theorem}

\begin{proof}
It remains to treat $r\to\infty$, since bounded subsequences follow
from Theorem~\ref{thm:fixed-precision-ensemble}. Split the canonical
cost as $K_\ell=R_{\ell,r}+B_{\ell,r}$, where $B_{\ell,r}$ is the sum
of the terminal $r$ folded costs. The Radon derivative of the law of
the terminal block conditional on $K_\ell=k$, relative to its
unconditioned law, is
\begin{equation}\label{eq:block-radon}
 D_{\ell,r}(B_{\ell,r})
 =\frac{\Pr(R_{\ell,r}=k-B_{\ell,r})}
        {\Pr(K_\ell=k)}.
\end{equation}
Choose $h_\ell\to\infty$ slowly enough that
$h_\ell\sqrt{r/\ell}\to0$. Uniform exponential moments give
\[
 \Pr\bigl(|B_{\ell,r}-\mathbb EB_{\ell,r}|>
                 h_\ell\sqrt r\bigr)=o(1).
\]
On the complementary event, removing the block changes the centered
local-CLT coordinate by $o(1)$ and changes the variance by a relative
$o(1)$. The uniform lattice local central limit theorem for the folded
geometric triangular array therefore gives
$D_{\ell,r}(B_{\ell,r})=1+o(1)$ uniformly there. Since
$\Pr(K_\ell=k)\gg_\delta\ell^{-1/2}$ in the central window, while the
standard concentration bound for lattice sums with uniformly bounded
exponential moments gives
$\sup_m\Pr(R_{\ell,r}=m)\ll(\ell-r)^{-1/2}$. Hence
$D_{\ell,r}\ll_\delta1$ on its whole support. Truncation of the
exceptional set now gives
$\mathbb E|D_{\ell,r}-1|=o(1)$. This is twice the total-variation
distance between the conditioned and unconditioned terminal blocks.
Their deterministic residue maps cannot increase that distance. The
cap-reduction identity in the proof of
Theorem~\ref{thm:fixed-precision-ensemble} shows that the unconditioned
residue has law exactly $\mu_r$.
\end{proof}

\begin{theorem}[Linear-block nonequivalence]\label{thm:linear-block-nonequivalence}
Let $\mathbf B_{\ell,r}=(J_{\ell-r},\ldots,J_{\ell-1})$ be the
terminal block of folded costs and let
$S_{\ell,r}=\sum_{i=\ell-r}^{\ell-1}J_i$. Suppose that
\[
 \frac r\ell\longrightarrow\rho\in(0,1),\qquad
 \frac{k_\ell-\ell}{\sqrt\ell}\longrightarrow u.
\]
Then
\begin{align*}
 &d_{\rm TV}\bigl(
   \mathcal L(\mathbf B_{\ell,r}\mid K_\ell=k_\ell),
   \mathcal L(\mathbf B_{\ell,r})\bigr)\\
 &\quad=d_{\rm TV}\bigl(
   \mathcal L(S_{\ell,r}\mid K_\ell=k_\ell),
   \mathcal L(S_{\ell,r})\bigr)\\
 &\quad\longrightarrow
 d_{\rm TV}\bigl(
   \mathcal N(\rho u,2\rho(1-\rho)),
   \mathcal N(0,2\rho)\bigr)>0,
\end{align*}
where the second parameter of $\mathcal N$ is the variance. In fact,
the first equality holds without the $o(1)$ at every finite level.
For $u=0$, the limit equals
\[
 2\left[
 \Phi\left(\sqrt{\frac{-\log(1-\rho)}{\rho}}\right)
 -\Phi\left(\sqrt{\frac{-(1-\rho)\log(1-\rho)}{\rho}}\right)
 \right].
\]
\end{theorem}

\begin{proof}
Write $R_{\ell,r}=K_\ell-S_{\ell,r}$. Relative to the
unconditioned block law, the conditioned law has density
\[
 \frac{\Pr(R_{\ell,r}=k_\ell-S_{\ell,r})}
      {\Pr(K_\ell=k_\ell)}.
\]
This density is a function of $S_{\ell,r}$ alone. Taking its absolute
deviation from one proves the finite-level equality of the two total
variation distances.

The folded costs have uniformly bounded exponential moments. Their
means and variances differ from $1$ and $2$, respectively, by summable
errors. The lattice local central limit theorem, applied separately to
$S_{\ell,r}$ and $R_{\ell,r}$, gives the unconditioned limit
$\mathcal N(0,2\rho)$ for
$(S_{\ell,r}-r)/\sqrt\ell$. After conditioning, its limiting density is
\[
 x\longmapsto
 \frac{\phi_{0,2\rho}(x)\phi_{0,2(1-\rho)}(u-x)}
      {\phi_{0,2}(u)}
 =\phi_{\rho u,2\rho(1-\rho)}(x).
\]
Uniform exponential tails promote the local convergence to convergence
in total variation. The two limiting variances are different, so the
distance is positive. When $u=0$, the two centered Gaussian densities
cross at
$|x|=\sqrt{-2(1-\rho)\log(1-\rho)}$; integration on the central
interval gives the displayed formula.
\end{proof}

Theorem~\ref{thm:linear-block-nonequivalence} shows that the condition
$r=o(\ell)$ in Theorem~\ref{thm:sublinear-precision-ensemble} is sharp
for the latent cost vectors. It does not prove nonequivalence after
projection modulo $3^r$: the residue map can discard the block sum,
and total variation can decrease under that projection. The
linear-precision residue and Fourier problems remain open.

\begin{theorem}[Full-precision likelihood domination]
\label{thm:ensemble-divergence}
Let $p_{\ell,k}^{(r)}$ be the projection of $p_{\ell,k}$ modulo $3^r$,
where $1\le r\le\ell$. For every set
$E\subseteq\mathbb Z/3^r\mathbb Z$,
\begin{equation}\label{eq:ensemble-domination}
 p_{\ell,k}^{(r)}(E)
 \le \Pr(K_\ell=k)^{-1}\mu_r(E).
\end{equation}
Consequently
\[
 D_{\rm KL}(p_{\ell,k}^{(r)}\Vert\mu_r)
 \le D_\infty(p_{\ell,k}^{(r)}\Vert\mu_r)
 \le-\log\Pr(K_\ell=k).
\]
Uniformly for $|k-\ell|\le\delta\sqrt\ell$, the last quantity is
$\frac12\log\ell+O_\delta(1)$.
\end{theorem}

\begin{proof}
On the canonical folded-cost space, the microcanonical law is the
conditional law given $K_\ell=k$. Hence, for the inverse image
$\widetilde E$ of $E$ under the residue map,
\[
 \Pr(\widetilde E\mid K_\ell=k)
 \le\frac{\Pr(\widetilde E)}{\Pr(K_\ell=k)}.
\]
The unconditioned residue modulo $3^r$ has law $\mu_r$, by the same cap
reduction used in Theorem~\ref{thm:fixed-precision-ensemble}. This proves
\eqref{eq:ensemble-domination}. It is the asserted bound on the
essential likelihood ratio, and relative entropy is bounded by its
logarithm. The lattice local central limit theorem gives
$\Pr(K_\ell=k)\asymp_\delta\ell^{-1/2}$ in the central window.
\end{proof}

Theorem~\ref{thm:ensemble-divergence} applies at full precision, but its
direction is insufficient for condition $(\star3)$. It proves that the
relative entropy per digit vanishes and rules out an exponential
likelihood separation from the canonical ensemble. Condition $(\star3)$
asks for a lower bound on every residue, so the unresolved part is the
lower tail of the likelihood ratio.

For $r=o(\ell)$ and $\xi_0\in\mathbb Z/3^r\mathbb Z$,
Theorem~\ref{thm:sublinear-precision-ensemble} gives
\begin{equation}\label{eq:fixed-conductor-fourier}
 \widehat p_{\ell,k}(3^{\ell-r}\xi_0)
 =\widehat\mu_r(\xi_0)+o(1),
\end{equation}
uniformly in $\xi_0$ and in the central cost window.
Thus the fixed-cost law does not approach Haar at fixed precision.
Fixed-conductor modes may persist. For example, $\mu_1$ assigns masses
$1/3$ and $2/3$ to the two units modulo $3$, so its nontrivial Fourier
coefficient has modulus $1/\sqrt3$. Any Fourier estimate for condition
$(\star3)$ must separate these exact coarse modes and control frequencies
whose conductor is comparable with $\ell$.

\begin{empirical}[Fixed-precision projection test]
\label{thm:fixed-precision-finite}
We projected the exact fixed-cost tables onto $3^r$, for $r=1,2,3$,
and compared them in total variation with $\mu_r$ obtained from the
independent Syracuse recursion. At $\ell=12$ and $k=\ell$, the
distances were $0.02050$, $0.04024$, and $0.06228$. At
$k=\ell+\lfloor\sqrt\ell\rfloor$, they were $0.04310$, $0.07902$,
and $0.12677$. At fixed $r$ and $k=\ell$, the distance decreases
monotonically over every level tested, $\ell=3,\ldots,13$. This
finite test checks the implementation and is not used in the proof.
At full precision and $\ell=13$, the distances were $0.25195$ at
$k=\ell$ and $0.26078$ at
$k=\ell+\lfloor\sqrt\ell\rfloor$. The corresponding precision-one
values were $0.01885$ and $0.04003$. These full-precision values do not
determine an asymptotic limit.
\end{empirical}

\begin{empirical}[Finite microcanonical coverage]\label{thm:microcanonical-finite}
We computed the multiplicities $g_\ell(k,a)$ exactly through
$\ell=12$, checking at every level that their residue sum equals the
independent bounded-composition count. At $\ell=12$, the fractions of
unit residues reached at
$k=\ell-\lfloor\sqrt\ell\rfloor,\ell,
\ell+\lfloor\sqrt\ell\rfloor$ are respectively
$0.2735$, $0.8074$, and $0.9973$. The minimum multiplicity is still
zero in all three cases. A separate Boolean recursion extends the
support calculation through $\ell=16$. At that level the central and
upper-window coverage fractions are $0.9194$ and $0.99999972$, while
the first fixed cost covering every residue equals $\ell+5$ for each
$10\le\ell\le16$. These values neither prove nor refute the uniform
asymptotic lower bound.
\end{empirical}

\begin{empirical}[Fourier budget for fixed-cost multiplicities]
\label{thm:microcanonical-fourier}
For $p_{\ell,k}(a)=g_\ell(k,a)/\sum_b g_\ell(k,b)$, additive Fourier
inversion gives
\[
 \min_a\frac{p_{\ell,k}(a)}{1/(2\cdot3^{\ell-1})}
 \ge 1-\frac23\sum_\xi
 |\widehat p_{\ell,k}(\xi)-\widehat u_\ell(\xi)|.
\]
At the central cost $k=\ell$, the spectral sum grows from $12.11$ at
$\ell=4$ to $293.29$ at $\ell=12$ and is dominated by primitive
frequencies. At $k=\ell+\lfloor\sqrt\ell\rfloor$, it grows from $7.61$
to $122.63$. At $\ell=12$, the largest individual discrepancy is
$0.339$ at central cost and $0.280$ in the upper window, while the
largest primitive discrepancies are $0.0246$ and $0.00682$. Thus the
largest coefficients retain a coarse component even as each primitive
coefficient becomes smaller. The resulting lower bounds are negative
throughout. This does not disprove uniform positivity; it shows that
termwise absolute summation of the finite Fourier series does not
establish it on this range.
\end{empirical}

\section{Certified numerical test of Wirsching's Conjecture 3}
\label{sec:conj3}

Conjecture 3 asserts $(\star5)$: there are constants $c,\delta_5>0$
with $\lim_\ell\varphi(z_\ell)/\varphi_0(z_\ell)=c>0$ uniformly for
sequences $(z_\ell)\in\widetilde A_{\delta_5}$
\cite[\S7]{Wirsching2003}. The source proves $(\star5)\Rightarrow
(\star4)$ only for $\delta_5>\delta$, strictly larger than the class
tested by $(\star4)$ itself; we test $(\star5)$ below without
re-deriving that implication. For $(x_\ell)\in\widetilde A_\delta$,
write $x_\ell^+:=x_\ell+3^{-\ell-1}$ \cite[p.16]{Wirsching2003}. The
test below uses $x_\ell=k_\ell/3^\ell$ with $k_\ell=\ell+\lfloor
u\sqrt\ell\rceil$ for $u\in\{-2,-1,-0.5,0,0.5,1,2\}$, the class
$\delta\approx2$ window announced as ``$u\in[-2,2]$'' in
Empirical Result~\ref{thm:conjecture3} below.

Using the exact self-similarity $X\overset{d}{=}(2U+X)/3$, the moments
$M_i=\mathbb{E}[X^i]$ satisfy the exact rational recursion
\[
  M_i = \frac{1}{3^i-1}\sum_{k=1}^i \binom{i}{k}\frac{2^k}{k+1} M_{i-k},
\]
and, combined with an antiderivative reduction expressing $\varphi$ at
extreme-tail arguments in terms of exact binomial moment sums (rather
than iterating $W_3$, whose $L^1$ error certificate is too weak to
control pointwise values at the relevant scale), we obtain a
zero-truncation-error evaluation of $\varphi$ at depth $\ell$ up to
$\ell=500$ (working throughout in log-space with $60$-digit precision;
error certified $\le 10^{-8}$).

\begin{empirical}[Numerical test of Wirsching's Conjecture 3]
\label{thm:conjecture3}
The ratio $L_\ell := 3^{1-\ell}\varphi(3x_\ell^+)/\varphi(x_\ell^+)$,
measured for $\ell=250,\dots,500$ in the CLT window ($u\in[-2,2]$),
converges to the predicted value $2/3$ with deficit
$\ell\cdot(2/3-L_\ell) \to 0.580\pm 0.001$, a coefficient matching
the asymptotic $\varphi_0$ of Wirsching's equation (7.11)
\cite[eq.~(7.11)]{Wirsching2003}, itself credited there to Berg and
Krüppel \cite{BergKruppel1998}. The quantity $\ln(\varphi/\varphi_0)$
is increasing, concave, and uniform across $u\in[-2,2]$ (dispersion
$<10^{-4}$), consistent with a finite limit $L=-0.619\pm0.001$ from
the fit alone, dominated by a tail-shape systematic uncertainty of
$\pm 0.015$, giving $c=e^L\approx 0.538$ (range $0.53$--$0.55$), best
fit by a tail term $0.79/\sqrt\ell$; no log-periodic modulation is
detected. The $10^{-8}$ error bound above certifies the evaluation of
$\varphi$ itself, not this extrapolation, whose actual uncertainty is
the $\pm0.015$ figure. Conjecture~2 remains open;
Theorem~\ref{thm:wirsching-conj1} settles Conjecture~1 independently
of this computation.
\end{empirical}

\section{Discussion}
\label{sec:discussion}

The mass decomposition \eqref{eq:microcanonical} and the equivalence
and nonequivalence theorems above are exact; they bound how far
local-limit methods reach toward Conjecture 2, not the conjecture
itself. Theorem~\ref{thm:fixed-precision-ensemble} and
Theorem~\ref{thm:sublinear-precision-ensemble} give convergence at
fixed and sublinear precision; Theorem~\ref{thm:linear-block-nonequivalence}
shows the latent cost vectors provably fail to converge at linear
precision. Theorem~\ref{thm:ensemble-divergence} is a separate,
one-sided bound at full precision, $D_\infty\le\frac12\log\ell+O(1)$:
it does not explain the linear-precision failure, since it only
upper-bounds the likelihood-ratio gap, in the direction opposite to
the lower bound condition $(\star3)$ actually needs. Condition
$(\star4)$, the hypothesis Conjecture~2 poses toward $(\star3)$,
carries no information about the coordinate this machinery
controls (\S\ref{sec:chain}); closing the gap needs an argument on
$W_3$ itself, or a bridge between $(\star4)$ and $(\star3)$ uniform in
growing resolution, which neither of the two the source supplies
(\S\ref{sec:chain}) reaches.

\section*{Data Availability Statement}

Code and data reproducing every claim in this paper, organized by
section with a README per subfolder, are at
\url{https://github.com/faculdade/collatz-wirsching-2003}.

\section*{Acknowledgments}

AI assistance (Claude, Anthropic) was used for textual review and
translation.

\bibliographystyle{plain}
\begin{thebibliography}{9}

\bibitem{Wirsching2003} G.\ J.\ Wirsching, \emph{On the problem of
  positive predecessor density in $3n+1$ dynamics}, Discrete and
  Continuous Dynamical Systems \textbf{9}(3) (2003), 771--787.
  DOI: 10.3934/dcds.2003.9.771.

\bibitem{BergKruppel1998} L.\ Berg and M.\ Kr\"uppel, \emph{On the
  solution of an integral-functional equation with a parameter},
  Zeitschrift f\"ur Analysis und ihre Anwendungen \textbf{17} (1998),
  159--181.

\bibitem{Wirsching1998Urns} G.\ J.\ Wirsching, \emph{Balls in
  constrained urns and Cantor-like sets}, Zeitschrift f\"ur Analysis
  und ihre Anwendungen \textbf{17} (1998), 979--996.

\bibitem{DragicevicEtAl2018} D.\ Dragi\v{c}evi\'c, G.\ Froyland,
  C.\ Gonz\'alez-Tokman, and S.\ Vaienti, \emph{A spectral approach
  for quenched limit theorems for random expanding dynamical systems},
  Communications in Mathematical Physics \textbf{360}(3) (2018),
  1121--1187. DOI: 10.1007/s00220-017-3083-7.

\bibitem{Hafouta2020} Y.\ Hafouta, \emph{Limit theorems for some
  time-dependent expanding dynamical systems}, Nonlinearity
  \textbf{33} (2020), 6421--6460.
  DOI: 10.1088/1361-6544/aba5e7.

\bibitem{BarrierCompanion} R.\ A.\ Tavares, \emph{The $qx+1$
  generalization of Tao's Syracuse measure: an endogeny barrier},
  companion paper, in preparation. Code and data:
  \url{https://github.com/faculdade/collatz-endogeny}.

\end{thebibliography}

\end{document}
